% Avsnitt 3.11, ur lectures/L03/appendix/b_exercises.md.

\section{Övningar}
\label{c3:sec:exercises}

\begin{aside}
\textbf{Så kontrollerar du ditt arbete.} Varje övning nedan som ber dig skriva en VHDL-modul har
en utdelad självkontrollerande testbänk under \file{lectures/L03/exercises/}. Skriv din modul i
dess katalog \file{exercises/<module>/}, med det entitetsnamn och den \textbf{portordning}
övningen anger, och kör den sedan med GHDL \textemdash{} se \secref{c2:sec:testbenches} för de tre
kommandona. Det här är det första kapitlet vars övningar behöver det för en klockad design, så det
är värt att läsa om innan du börjar med \code{register4}.
\end{aside}

% ------------------------------------------------------------------------------------------
\exerciseset{D-låset}

\exercise{Den första kretsen som minns}{Krets}
\label{c3:ex:latch}
Bygg den första kretsen i den här boken som kan \emph{minnas} något.

Allt du har byggt hittills har varit kombinatoriskt: utgångarna följer ingångarna, och i samma
ögonblick som ingångarna ändras ändras utgångarna med dem. Ett lås är det minsta möjliga avsteget
från det. Det har en dataingång \code{D} och en styringång \code{enable}, och det har två
utgångar, \code{Q} och dess komplement \code{Qn}. Så länge \code{enable} är hög är låset
\textbf{genomsläppligt}: \code{Q} följer helt enkelt \code{D}, och det ser ut som en ledning. Det
intressanta är vad som händer när \code{enable} går låg. Låset \textbf{låser}, och håller kvar det
\code{D} var i det ögonblicket, och från och med då kan \code{D} göra vad det vill utan att
\code{Q} märker det. Det hållna värdet är en bit minne, och det är byggt av ingenting annat än de
grindar du redan känner till, kopplade så att utgången återkopplas till ingången.

\code{Qn} är ingen dekoration. Återkopplingsslingan är det som får kretsen att minnas, och de två
utgångarna är det som matar varandra: \code{Q} beräknas ur \code{Qn} och \code{Qn} ur \code{Q}.
Bygg det så ser du slingan; det är hela poängen med att rita det för hand innan du ens möter
VHDL:en.

Använd ekvationerna från \secref{c3:sec:latch}:

\[ Q = (D' \cdot \mathit{enable} + Q_n)' \qquad Q_n = (D \cdot \mathit{enable} + Q)' \]

\xpart{a} Realisera motsvarande grindnät: rita nätet för hand, och återskapa och simulera det i
CircuitVerse.

\xpart{b} Testa det genomsläppliga och det låsta tillståndet. Sätt \code{enable = 1}, ändra
\code{D}, och observera vad som händer med \code{Q} och \code{Qn}. Sätt sedan \code{enable = 0},
ändra \code{D} igen, och observera vad som händer. Förklara skillnaden mellan de två fallen.

\xpart{c} Kör \code{D} och \code{enable} genom kombinationerna \code{00}, \code{01}, \code{10},
\code{11}, två gånger. Bekräfta att du observerar båda fallen: låset låser efter att ha varit
genomsläppligt med \code{D = 1}, och låset låser efter att ha varit genomsläppligt med
\code{D = 0}. Förklara hur värdet på \code{Q} beror på värdet på \code{D} i det ögonblick
\code{enable} ändras från \code{1} till \code{0}.

% ------------------------------------------------------------------------------------------
\exerciseset{D-vippan}

\exercise{Från lås till vippa}{Krets}
\label{c3:ex:dff}
Bygg ut D-låset från Övning~\ref{c3:ex:latch} till en D-vippa, enligt \secref{c3:sec:dff}.

Vippan ska ha ingångarna \code{clock}, \code{reset_n}, \code{D} och \code{enable}, och utgångarna
\code{Q} och \code{Qn}.

Resetsignalen är \textbf{asynkron} (den får verkan utan att vänta på en klockflank) och
\textbf{aktiv låg} (att sätta \code{reset_n = 0} ska tvinga \code{Q} till \code{0} omedelbart).

\task{Tips}
\begin{itemize}
  \item Resetten måste göra två saker i vart och ett av de två interna låsen, inte en:
  \begin{itemize}
    \item Mata in ett aktivt högt \code{reset} (\code{reset_n} genom en NOT-grind) i den NOR-grind
      som producerar \code{Q}. Det tvingar \code{Q} låg utan att vänta på en klockflank.
    \item \code{AND}:a in \code{reset_n} i den AND-grind som producerar det låsets produkt av
      \code{D} och \code{enable}. Det frigör \code{Qn} att gå hög.
    \item Att grinda enbart återkopplingsvägen för \code{Q} räcker inte: med \code{enable = 1} och
      \code{D = 1} lämnar det \code{Qn} låg, och en låg \code{Qn} är precis det som håller \code{Q}
      hög, så resetten skulle inte göra någonting.
  \end{itemize}
  \item Driv de två interna låsens enable-ingångar med \code{clock} respektive inversen av
    \code{clock}.
  \item Följ \secref{c3:sec:dff} för hur den yttre \code{enable}-ingången styr om vippan tar emot
    ett nytt värde.
\end{itemize}

\xpart{a} Realisera motsvarande grindnät: rita det för hand, återskapa och simulera det i
CircuitVerse, och sätt klockperioden till \code{1000 ms}, långsamt nog för att observera beteendet
med ögat.

\xpart{b} Testa det klockade beteendet: ändra \code{D} under klockans höga fas, och sedan under
dess låga fas, och observera \code{Q} och \code{Qn}. Bekräfta att utgångarna ändras bara vid den
aktiva klockflanken, aldrig mellan klockflankerna. Förklara hur de två interna låsen samverkar för
att ge flanktriggat beteende.

\xpart{c} Testa den asynkrona resetten: sätt först \code{Q = 1}, aktivera sedan resetten genom att
sätta \code{reset_n = 0}, och observera om \code{Q} nollställs omedelbart eller vid nästa
klockflank. Förklara resultatet genom att hänvisa till skillnaden mellan en asynkron signal och en
synkron signal.

% ------------------------------------------------------------------------------------------
\exerciseset{Register}

\exercise{\code{register4}}{VHDL}
\label{c3:ex:register4}
Skriv en fullständig VHDL-entitet med namnet \code{register4}: fyra vippor sida vid sida, som delar
en klocka, en reset och en enable.

Det här är övningen där vippan slutar vara en kuriositet och blir en byggsten. En ensam D-vippa
lagrar en bit, vilket inte är till mycket nytta i sig; vad en design egentligen vill ha är
någonstans att förvara ett \emph{tal}. Sätt fyra vippor på rad, koppla samma klocka till dem alla,
och du har ett fyrabitars register. Ingenting hos den enskilda vippan ändras. Bara bredden gör det.

Det som är värt att lägga märke till är hur lite VHDL:en behöver ändras för att säga det. Du
skriver samma synkrona processmall från \secref{c3:sec:template} som du redan använde för en bit,
med \code{std_logic_vector(3 downto 0)} i stället för \code{std_logic}. Det finns ingen loop, ingen
upprepning, inga fyra kopior av någonting. Det är ingen genväg som språket erbjuder dig; det är ett
faktum om hårdvaran, och del \textbf{c)} nedan ber dig säga varför.

Entiteten ska ha:

\begin{booktable}{@{}p{5em}p{3.4em}p{10.8em}L@{}}
  \thead{Port} & \thead{Riktn.} & \thead{Typ} & \thead{Beskrivning} \\
  \midrule
  \code{clock} & in & \code{std_logic} & Systemklocka. \\
  \code{reset_n} & in & \code{std_logic} & Asynkron, aktiv låg. Nollställer registret till
    \code{"0000"}. \\
  \code{enable} & in & \code{std_logic} & \code{'0'} behåller nuvarande värde; \code{'1'} fångar
    \code{d} vid nästa stigande flank. \\
  \code{d} & in & \code{std_logic_vector(3 downto 0)} & Värdet som ska fångas. \\
  \code{q} & out & \code{std_logic_vector(3 downto 0)} & Det lagrade värdet. \\
\end{booktable}

Skriv ut processen själv i stället för att kopiera det genomarbetade vippexemplet i
\secref{c3:sec:dffvhdl}. Att knappa in den är vad som får mallen att sluta vara något du känner
igen och börja vara något du kan producera ur minnet, och du kommer att skriva den i varenda
återstående kapitel i den här boken.

\xpart{a} Skriv entiteten och den synkrona processen, och anpassa mallen från en enda bit till en
fyrabitars vektor.

\xpart{b} Kontrollera den med testbänken, så som beskrivs högst upp i det här avsnittet.

\xpart{c} Förklara varför samma mall skalar från en vippa till fyra utan någon strukturell ändring,
medan en fyrabitars \emph{räknare} inte helt enkelt kan vara fyra enbitsräknare sida vid sida. Vad
skiljer de två fallen åt?

\modulefig[0.6]{lectures/L03/appendix/images/register4.png}

\begin{selfcheck}
\textbf{Självkontroll.} Döp din entitet till \code{register4}, med portarna \code{clock},
\code{reset_n}, \code{enable}, \code{d} (\code{std_logic_vector(3 downto 0)}) och \code{q}
(\code{std_logic_vector(3 downto 0)}), deklarerade i den ordningen; dess testbänk finns i
\file{exercises/register4/}.
\end{selfcheck}

% ------------------------------------------------------------------------------------------
\exerciseset{Flankdetektering}

\exercise{Växlingskretsen för hand}{Krets}
\label{c3:ex:edgehand}
Konstruera en krets som växlar en lysdiod en gång för varje samplad stigande flank hos en
knappsignal. Följ \secref{c3:sec:edge} och \secref{c3:sec:ledtoggle}.

Kretsen ska ha ingångarna \code{button}, \code{clock} och \code{reset}, och utgången \code{led}.

I den här övningen är \code{reset} aktiv hög, till skillnad från resetten i det genomarbetade
exemplet. Den är fortfarande \textbf{asynkron}, som varje reset i den här boken: aktivera den så
nollställs \code{led} omedelbart, utan att vänta på en klockflank. I mallen från
\secref{c3:sec:template} betyder det att resetgrenen stannar utanför testet
\code{rising_edge(clock)}, och att bara polariteten på resettestet ändras.

\task{Tips}
Använd två vippor. Den ena lagrar det föregående samplade värdet av \code{button}, och det värdet
används för att detektera en stigande flank. Den andra lagrar lysdiodens nuvarande tillstånd, och
växlas bara när en stigande flank detekteras.

\xpart{a} Realisera motsvarande grindnät: rita det för hand, återskapa och simulera det i
CircuitVerse, och sätt klockperioden till \code{1000 ms}.

\xpart{b} Testa flankdetektorn: håll \code{button} hög under flera klockcykler i följd, bekräfta
att \code{led} växlar exakt en gång, och att den inte växlar en gång per klockcykel. Förklara
vilken lagrad signal som hindrar kretsen från att gång på gång detektera samma höga nivå som en ny
stigande flank.

\exercise{\code{led_toggle_single}}{VHDL}
\label{c3:ex:ledsingle}
Implementera kretsen från Övning~\ref{c3:ex:edgehand} i VHDL.

\xpart{a} Skapa en entitet med namnet \code{led_toggle_single} med:

\begin{booktable}{@{}p{5em}p{3.4em}p{5.6em}L@{}}
  \thead{Port} & \thead{Riktn.} & \thead{Typ} & \thead{Beskrivning} \\
  \midrule
  \code{clock} & in & \code{std_logic} & Systemklocka. \\
  \code{reset} & in & \code{std_logic} & Asynkron, och \textbf{aktiv hög}, till skillnad från varje
    annan modul i den här boken. Aktivera den så nollställs \code{led} utan att vänta på en
    klockflank, så resetgrenen hör hemma utanför testet \code{rising_edge(clock)}. \\
  \code{button} & in & \code{std_logic} & Växlar \code{led} en gång per tryck. \\
  \code{led} & out & \code{std_logic} & Lysdiodens tillstånd. \\
\end{booktable}

Lägg märke till enbitstyperna: det genomarbetade exemplet använder
\code{std_logic_vector}-portar för två knappar och två lysdioder, medan den här driver ett enda
par.

Kopiera inte \file{led_toggle.vhd} rakt av. Gör så här i stället:
\begin{itemize}
  \item Härled de två klockade processerna själv: den ena lagrar det föregående värdet av
    \code{button}, den andra lagrar och växlar det nuvarande värdet av \code{led}.
  \item Härled flankdetekteringsekvationen för en stigande flank och en aktivt hög knapp, alltså
    \secref{c3:sec:edge}:s $\mathit{edge} = \mathit{current} \cdot \mathit{previous}'$.
  \item Använd inte ekvationen för fallande flank från det genomarbetade exemplet, som utgår från
    en aktivt låg knapp.
\end{itemize}

\xpart{b} Verifiera designen med dess testbänk.

Testbänken driver samma sekvens som du skulle gå igenom för hand på ett kort:
\begin{itemize}
  \item Den aktiverar \code{reset} och kontrollerar att \code{led} är nollställd.
  \item Den drar \code{button} hög och kontrollerar att \code{led} växlar exakt en gång.
  \item Den \textbf{håller \code{button} hög} under ytterligare några klockcykler och kontrollerar,
    efter var och en av dem, att \code{led} \emph{inte} växlar igen.
  \item Den släpper och trycker igen, och kontrollerar att det andra trycket växlar \code{led} en
    gång till.
  \item Den aktiverar \code{reset} med \code{led} tänd och \textbf{utan någon klockflank efteråt},
    och kontrollerar att \code{led} nollställs ändå. Det är den kontroll som faller om du lägger
    resetten inuti grenen \code{rising_edge(clock)}.
\end{itemize}

Den tredje kontrollen är den som faller om din flankdetektor är fel. En \code{led} som växlar varje
klockcykel medan knappen hålls nedtryckt betyder att du drev växlingen från \code{button}s
\textbf{nivå} i stället för från dess flank. Läs om ekvationen ovan: \code{edge} är hög bara när
\code{button} är hög \emph{nu} och dess lagrade föregående värde är lågt.

\modulefig[0.56]{lectures/L03/appendix/images/led_toggle_single.png}

\begin{selfcheck}
\textbf{Självkontroll.} Döp din entitet till \code{led_toggle_single}, med portarna \code{clock},
\code{reset}, \code{button} (in) och \code{led} (out), deklarerade i den ordningen; dess testbänk
finns i \file{exercises/led_toggle_single/}. Lägg märke till att \code{reset} är \textbf{aktiv hög}
här, till skillnad från i det genomarbetade exemplet.
\end{selfcheck}
